\documentclass[12pt,letterpaper]{article}
\usepackage[margin=0in]{geometry}
\usepackage{tikz}
\usepackage[T1]{fontenc}
\usepackage{lmodern}\pdfmapfile{+lm.map}
\pagestyle{empty}
\setlength{\parindent}{0pt}\hyphenpenalty=10000\exhyphenpenalty=10000
\begin{document}
\null
\begin{tikzpicture}[remember picture,overlay,x=1mm,y=-1mm]
\begin{scope}[shift={(current page.north west)}]
\node[anchor=north west,inner sep=0pt,outer sep=0pt,align=left,text width=187mm] at (15,12) {\fontsize{11}{13.86}\selectfont Week 26 / Same area, different boundaries / Grades 4--5};
\node[anchor=north west,inner sep=0pt,outer sep=0pt,align=left,text width=184mm] at (15,26) {\fontsize{12}{15.12}\selectfont Keep all tiles flat and in one piece, joined along whole sides without overlaps. Count every side with no tile beside it, including sides around holes. Turns and flips count as the same shape.};
\node[anchor=north west,inner sep=0pt,outer sep=0pt,align=left,text width=184mm] at (15,66) {\fontsize{14}{17.64}\selectfont \textbf{Problem 1:} Find the shortest boundary you can with 7, 10, and 13 tiles. For each number, find two different shapes reaching your best length, if you can.};
\begin{scope}[shift={(20,103)}]
\draw[step=8mm,black!27,line width=.3pt] (0,0) grid (56,40);
\end{scope}
\draw[black!27,line width=.3pt] (20,103) rectangle (76,143);
\node[anchor=north west,inner sep=0pt,outer sep=0pt,align=left] at (20,145) {\fontsize{11}{13.86}\selectfont 7 tiles; boundary: \rule{10mm}{.3pt}};
\begin{scope}[shift={(119,103)}]
\draw[step=8mm,black!27,line width=.3pt] (0,0) grid (56,40);
\end{scope}
\draw[black!27,line width=.3pt] (119,103) rectangle (175,143);
\node[anchor=north west,inner sep=0pt,outer sep=0pt,align=left] at (119,145) {\fontsize{11}{13.86}\selectfont 7 tiles; boundary: \rule{10mm}{.3pt}};
\begin{scope}[shift={(20,155)}]
\draw[step=8mm,black!27,line width=.3pt] (0,0) grid (56,40);
\end{scope}
\draw[black!27,line width=.3pt] (20,155) rectangle (76,195);
\node[anchor=north west,inner sep=0pt,outer sep=0pt,align=left] at (20,197) {\fontsize{11}{13.86}\selectfont 10 tiles; boundary: \rule{10mm}{.3pt}};
\begin{scope}[shift={(119,155)}]
\draw[step=8mm,black!27,line width=.3pt] (0,0) grid (56,40);
\end{scope}
\draw[black!27,line width=.3pt] (119,155) rectangle (175,195);
\node[anchor=north west,inner sep=0pt,outer sep=0pt,align=left] at (119,197) {\fontsize{11}{13.86}\selectfont 10 tiles; boundary: \rule{10mm}{.3pt}};
\begin{scope}[shift={(20,207)}]
\draw[step=8mm,black!27,line width=.3pt] (0,0) grid (56,40);
\end{scope}
\draw[black!27,line width=.3pt] (20,207) rectangle (76,247);
\node[anchor=north west,inner sep=0pt,outer sep=0pt,align=left] at (20,249) {\fontsize{11}{13.86}\selectfont 13 tiles; boundary: \rule{10mm}{.3pt}};
\begin{scope}[shift={(119,207)}]
\draw[step=8mm,black!27,line width=.3pt] (0,0) grid (56,40);
\end{scope}
\draw[black!27,line width=.3pt] (119,207) rectangle (175,247);
\node[anchor=north west,inner sep=0pt,outer sep=0pt,align=left] at (119,249) {\fontsize{11}{13.86}\selectfont 13 tiles; boundary: \rule{10mm}{.3pt}};
\node[anchor=north west,inner sep=0pt,outer sep=0pt,align=left] at (15,268) {\fontsize{9}{11.34}\selectfont Bellingham Math Circle / Week 26 / W26-45-v2};
\node[anchor=north east,inner sep=0pt,outer sep=0pt,align=left] at (199,268) {\fontsize{9}{11.34}\selectfont 1};
\end{scope}
\end{tikzpicture}
\newpage
\null
\begin{tikzpicture}[remember picture,overlay,x=1mm,y=-1mm]
\begin{scope}[shift={(current page.north west)}]
\node[anchor=north west,inner sep=0pt,outer sep=0pt,align=left,text width=187mm] at (15,12) {\fontsize{11}{13.86}\selectfont Week 26 / Same area, different boundaries / Grades 4--5};
\node[anchor=north west,inner sep=0pt,outer sep=0pt,align=left,text width=184mm] at (15,28) {\fontsize{14}{17.64}\selectfont \textbf{Problem 2:} What is the longest possible boundary of a twelve-tile shape? Explain why no shape can have a longer boundary. Draw three shapes reaching that length, with none in a straight row. Can such a shape contain a 2-by-2 block? Can it have a hole?};
\begin{scope}[shift={(15,78)}]
\draw[step=5.5mm,black!27,line width=.3pt] (0,0) grid (77.0,27.5);
\end{scope}
\draw[black!27,line width=.3pt] (15,78) rectangle (92.0,105.5);
\node[anchor=north west,inner sep=0pt,outer sep=0pt,align=left] at (15,109) {\fontsize{11}{13.86}\selectfont Boundary: \rule{14mm}{.3pt}};
\begin{scope}[shift={(114,78)}]
\draw[step=5.5mm,black!27,line width=.3pt] (0,0) grid (77.0,27.5);
\end{scope}
\draw[black!27,line width=.3pt] (114,78) rectangle (191.0,105.5);
\node[anchor=north west,inner sep=0pt,outer sep=0pt,align=left] at (114,109) {\fontsize{11}{13.86}\selectfont Boundary: \rule{14mm}{.3pt}};
\begin{scope}[shift={(15,138)}]
\draw[step=5.5mm,black!27,line width=.3pt] (0,0) grid (77.0,27.5);
\end{scope}
\draw[black!27,line width=.3pt] (15,138) rectangle (92.0,165.5);
\node[anchor=north west,inner sep=0pt,outer sep=0pt,align=left] at (15,169) {\fontsize{11}{13.86}\selectfont Boundary: \rule{14mm}{.3pt}};
\draw[black!30,line width=.35pt] (15,184) -- ++(179,0);
\draw[black!30,line width=.35pt] (15,202) -- ++(179,0);
\draw[black!30,line width=.35pt] (15,220) -- ++(179,0);
\draw[black!30,line width=.35pt] (15,238) -- ++(179,0);
\draw[black!30,line width=.35pt] (15,256) -- ++(179,0);
\node[anchor=north west,inner sep=0pt,outer sep=0pt,align=left] at (15,268) {\fontsize{9}{11.34}\selectfont Bellingham Math Circle / Week 26 / W26-45-v2};
\node[anchor=north east,inner sep=0pt,outer sep=0pt,align=left] at (199,268) {\fontsize{9}{11.34}\selectfont 2};
\end{scope}
\end{tikzpicture}
\newpage
\null
\begin{tikzpicture}[remember picture,overlay,x=1mm,y=-1mm]
\begin{scope}[shift={(current page.north west)}]
\node[anchor=north west,inner sep=0pt,outer sep=0pt,align=left,text width=187mm] at (15,12) {\fontsize{11}{13.86}\selectfont Week 26 / Same area, different boundaries / Grades 4--5};
\node[anchor=north west,inner sep=0pt,outer sep=0pt,align=left,text width=184mm] at (15,28) {\fontsize{14}{17.64}\selectfont \textbf{Problem 3:} Count the occupied rows, occupied columns, and boundary sides of each shape.};
\filldraw[fill=black!13,draw=black,line width=.75pt] (18,67) rectangle ++(7,7);
\filldraw[fill=black!13,draw=black,line width=.75pt] (18,74) rectangle ++(7,7);
\filldraw[fill=black!13,draw=black,line width=.75pt] (25,67) rectangle ++(7,7);
\filldraw[fill=black!13,draw=black,line width=.75pt] (25,74) rectangle ++(7,7);
\filldraw[fill=black!13,draw=black,line width=.75pt] (32,67) rectangle ++(7,7);
\filldraw[fill=black!13,draw=black,line width=.75pt] (32,74) rectangle ++(7,7);
\filldraw[fill=black!13,draw=black,line width=.75pt] (39,67) rectangle ++(7,7);
\filldraw[fill=black!13,draw=black,line width=.75pt] (39,74) rectangle ++(7,7);
\filldraw[fill=black!13,draw=black,line width=.75pt] (46,67) rectangle ++(7,7);
\filldraw[fill=black!13,draw=black,line width=.75pt] (46,74) rectangle ++(7,7);
\filldraw[fill=black!13,draw=black,line width=.75pt] (53,67) rectangle ++(7,7);
\filldraw[fill=black!13,draw=black,line width=.75pt] (53,74) rectangle ++(7,7);
\node[anchor=north west,inner sep=0pt,outer sep=0pt,align=left] at (18,99) {\fontsize{11}{13.86}\selectfont Rows: \rule{8mm}{.3pt}};
\node[anchor=north west,inner sep=0pt,outer sep=0pt,align=left] at (18,106) {\fontsize{11}{13.86}\selectfont Columns: \rule{8mm}{.3pt}};
\node[anchor=north west,inner sep=0pt,outer sep=0pt,align=left] at (18,113) {\fontsize{11}{13.86}\selectfont Boundary: \rule{8mm}{.3pt}};
\filldraw[fill=black!13,draw=black,line width=.75pt] (85,67) rectangle ++(7,7);
\filldraw[fill=black!13,draw=black,line width=.75pt] (85,74) rectangle ++(7,7);
\filldraw[fill=black!13,draw=black,line width=.75pt] (85,81) rectangle ++(7,7);
\filldraw[fill=black!13,draw=black,line width=.75pt] (92,67) rectangle ++(7,7);
\filldraw[fill=black!13,draw=black,line width=.75pt] (92,74) rectangle ++(7,7);
\filldraw[fill=black!13,draw=black,line width=.75pt] (92,81) rectangle ++(7,7);
\filldraw[fill=black!13,draw=black,line width=.75pt] (99,67) rectangle ++(7,7);
\filldraw[fill=black!13,draw=black,line width=.75pt] (99,74) rectangle ++(7,7);
\filldraw[fill=black!13,draw=black,line width=.75pt] (99,81) rectangle ++(7,7);
\filldraw[fill=black!13,draw=black,line width=.75pt] (106,67) rectangle ++(7,7);
\filldraw[fill=black!13,draw=black,line width=.75pt] (106,74) rectangle ++(7,7);
\filldraw[fill=black!13,draw=black,line width=.75pt] (106,81) rectangle ++(7,7);
\node[anchor=north west,inner sep=0pt,outer sep=0pt,align=left] at (85,99) {\fontsize{11}{13.86}\selectfont Rows: \rule{8mm}{.3pt}};
\node[anchor=north west,inner sep=0pt,outer sep=0pt,align=left] at (85,106) {\fontsize{11}{13.86}\selectfont Columns: \rule{8mm}{.3pt}};
\node[anchor=north west,inner sep=0pt,outer sep=0pt,align=left] at (85,113) {\fontsize{11}{13.86}\selectfont Boundary: \rule{8mm}{.3pt}};
\filldraw[fill=black!13,draw=black,line width=.75pt] (146,67) rectangle ++(7,7);
\filldraw[fill=black!13,draw=black,line width=.75pt] (153,67) rectangle ++(7,7);
\filldraw[fill=black!13,draw=black,line width=.75pt] (160,67) rectangle ++(7,7);
\filldraw[fill=black!13,draw=black,line width=.75pt] (167,67) rectangle ++(7,7);
\filldraw[fill=black!13,draw=black,line width=.75pt] (174,67) rectangle ++(7,7);
\filldraw[fill=black!13,draw=black,line width=.75pt] (146,74) rectangle ++(7,7);
\filldraw[fill=black!13,draw=black,line width=.75pt] (153,74) rectangle ++(7,7);
\filldraw[fill=black!13,draw=black,line width=.75pt] (160,74) rectangle ++(7,7);
\filldraw[fill=black!13,draw=black,line width=.75pt] (146,81) rectangle ++(7,7);
\filldraw[fill=black!13,draw=black,line width=.75pt] (153,81) rectangle ++(7,7);
\filldraw[fill=black!13,draw=black,line width=.75pt] (146,88) rectangle ++(7,7);
\filldraw[fill=black!13,draw=black,line width=.75pt] (153,88) rectangle ++(7,7);
\node[anchor=north west,inner sep=0pt,outer sep=0pt,align=left] at (146,99) {\fontsize{11}{13.86}\selectfont Rows: \rule{8mm}{.3pt}};
\node[anchor=north west,inner sep=0pt,outer sep=0pt,align=left] at (146,106) {\fontsize{11}{13.86}\selectfont Columns: \rule{8mm}{.3pt}};
\node[anchor=north west,inner sep=0pt,outer sep=0pt,align=left] at (146,113) {\fontsize{11}{13.86}\selectfont Boundary: \rule{8mm}{.3pt}};
\node[anchor=north west,inner sep=0pt,outer sep=0pt,align=left,text width=184mm] at (15,134) {\fontsize{14}{17.64}\selectfont Make two twelve-tile shapes that each occupy 4 rows and 5 columns but have different boundary lengths. What is the least possible boundary length for a shape occupying 4 rows and 5 columns? Explain.};
\begin{scope}[shift={(28,184)}]
\draw[step=11mm,black!27,line width=.3pt] (0,0) grid (55,44);
\end{scope}
\draw[black!27,line width=.3pt] (28,184) rectangle (83,228);
\node[anchor=north west,inner sep=0pt,outer sep=0pt,align=left] at (28,232) {\fontsize{11}{13.86}\selectfont Boundary: \rule{12mm}{.3pt}};
\begin{scope}[shift={(127,184)}]
\draw[step=11mm,black!27,line width=.3pt] (0,0) grid (55,44);
\end{scope}
\draw[black!27,line width=.3pt] (127,184) rectangle (182,228);
\node[anchor=north west,inner sep=0pt,outer sep=0pt,align=left] at (127,232) {\fontsize{11}{13.86}\selectfont Boundary: \rule{12mm}{.3pt}};
\draw[black!30,line width=.35pt] (15,247) -- ++(179,0);
\draw[black!30,line width=.35pt] (15,260) -- ++(179,0);
\node[anchor=north west,inner sep=0pt,outer sep=0pt,align=left] at (15,268) {\fontsize{9}{11.34}\selectfont Bellingham Math Circle / Week 26 / W26-45-v2};
\node[anchor=north east,inner sep=0pt,outer sep=0pt,align=left] at (199,268) {\fontsize{9}{11.34}\selectfont 3};
\end{scope}
\end{tikzpicture}
\newpage
\null
\begin{tikzpicture}[remember picture,overlay,x=1mm,y=-1mm]
\begin{scope}[shift={(current page.north west)}]
\node[anchor=north west,inner sep=0pt,outer sep=0pt,align=left,text width=187mm] at (15,12) {\fontsize{11}{13.86}\selectfont Week 26 / Same area, different boundaries / Grades 4--5};
\node[anchor=north west,inner sep=0pt,outer sep=0pt,align=left,text width=184mm] at (15,28) {\fontsize{14}{17.64}\selectfont \textbf{Problem 4:} For each boundary length shown, find the greatest number of tiles a shape can have. Explain why no shape with that boundary length can contain more tiles.};
\begin{scope}[shift={(24,70)}]
\draw[step=9mm,black!27,line width=.3pt] (0,0) grid (63,54);
\end{scope}
\draw[black!27,line width=.3pt] (24,70) rectangle (87,124);
\node[anchor=north west,inner sep=0pt,outer sep=0pt,align=left] at (24,127) {\fontsize{11}{13.86}\selectfont Boundary: 12; tiles: \rule{10mm}{.3pt}};
\begin{scope}[shift={(123,70)}]
\draw[step=9mm,black!27,line width=.3pt] (0,0) grid (63,54);
\end{scope}
\draw[black!27,line width=.3pt] (123,70) rectangle (186,124);
\node[anchor=north west,inner sep=0pt,outer sep=0pt,align=left] at (123,127) {\fontsize{11}{13.86}\selectfont Boundary: 14; tiles: \rule{10mm}{.3pt}};
\begin{scope}[shift={(24,156)}]
\draw[step=9mm,black!27,line width=.3pt] (0,0) grid (63,54);
\end{scope}
\draw[black!27,line width=.3pt] (24,156) rectangle (87,210);
\node[anchor=north west,inner sep=0pt,outer sep=0pt,align=left] at (24,213) {\fontsize{11}{13.86}\selectfont Boundary: 16; tiles: \rule{10mm}{.3pt}};
\begin{scope}[shift={(123,156)}]
\draw[step=9mm,black!27,line width=.3pt] (0,0) grid (63,54);
\end{scope}
\draw[black!27,line width=.3pt] (123,156) rectangle (186,210);
\node[anchor=north west,inner sep=0pt,outer sep=0pt,align=left] at (123,213) {\fontsize{11}{13.86}\selectfont Boundary: 18; tiles: \rule{10mm}{.3pt}};
\draw[black!30,line width=.35pt] (15,241) -- ++(179,0);
\draw[black!30,line width=.35pt] (15,257) -- ++(179,0);
\node[anchor=north west,inner sep=0pt,outer sep=0pt,align=left] at (15,268) {\fontsize{9}{11.34}\selectfont Bellingham Math Circle / Week 26 / W26-45-v2};
\node[anchor=north east,inner sep=0pt,outer sep=0pt,align=left] at (199,268) {\fontsize{9}{11.34}\selectfont 4};
\end{scope}
\end{tikzpicture}
\newpage
\null
\begin{tikzpicture}[remember picture,overlay,x=1mm,y=-1mm]
\begin{scope}[shift={(current page.north west)}]
\node[anchor=north west,inner sep=0pt,outer sep=0pt,align=left,text width=187mm] at (15,12) {\fontsize{11}{13.86}\selectfont Week 26 / Same area, different boundaries / Grades 4--5};
\node[anchor=north west,inner sep=0pt,outer sep=0pt,align=left,text width=184mm] at (15,28) {\fontsize{14}{17.64}\selectfont \textbf{Problem 5:} Find the shortest possible boundary for shapes with 12, 13, 17, 20, and 21 tiles. Explain why each of your shapes has the shortest possible boundary.};
\begin{scope}[shift={(24,66)}]
\draw[step=7mm,black!27,line width=.3pt] (0,0) grid (49,42);
\end{scope}
\draw[black!27,line width=.3pt] (24,66) rectangle (73,108);
\node[anchor=north west,inner sep=0pt,outer sep=0pt,align=left] at (24,111) {\fontsize{11}{13.86}\selectfont 12 tiles; boundary: \rule{10mm}{.3pt}};
\begin{scope}[shift={(123,66)}]
\draw[step=7mm,black!27,line width=.3pt] (0,0) grid (49,42);
\end{scope}
\draw[black!27,line width=.3pt] (123,66) rectangle (172,108);
\node[anchor=north west,inner sep=0pt,outer sep=0pt,align=left] at (123,111) {\fontsize{11}{13.86}\selectfont 13 tiles; boundary: \rule{10mm}{.3pt}};
\begin{scope}[shift={(24,126)}]
\draw[step=7mm,black!27,line width=.3pt] (0,0) grid (49,42);
\end{scope}
\draw[black!27,line width=.3pt] (24,126) rectangle (73,168);
\node[anchor=north west,inner sep=0pt,outer sep=0pt,align=left] at (24,171) {\fontsize{11}{13.86}\selectfont 17 tiles; boundary: \rule{10mm}{.3pt}};
\begin{scope}[shift={(123,126)}]
\draw[step=7mm,black!27,line width=.3pt] (0,0) grid (49,42);
\end{scope}
\draw[black!27,line width=.3pt] (123,126) rectangle (172,168);
\node[anchor=north west,inner sep=0pt,outer sep=0pt,align=left] at (123,171) {\fontsize{11}{13.86}\selectfont 20 tiles; boundary: \rule{10mm}{.3pt}};
\begin{scope}[shift={(24,186)}]
\draw[step=7mm,black!27,line width=.3pt] (0,0) grid (49,42);
\end{scope}
\draw[black!27,line width=.3pt] (24,186) rectangle (73,228);
\node[anchor=north west,inner sep=0pt,outer sep=0pt,align=left] at (24,231) {\fontsize{11}{13.86}\selectfont 21 tiles; boundary: \rule{10mm}{.3pt}};
\draw[black!30,line width=.35pt] (114,194) -- ++(80,0);
\draw[black!30,line width=.35pt] (114,210) -- ++(80,0);
\draw[black!30,line width=.35pt] (114,226) -- ++(80,0);
\draw[black!30,line width=.35pt] (15,245) -- ++(179,0);
\draw[black!30,line width=.35pt] (15,259) -- ++(179,0);
\node[anchor=north west,inner sep=0pt,outer sep=0pt,align=left] at (15,268) {\fontsize{9}{11.34}\selectfont Bellingham Math Circle / Week 26 / W26-45-v2};
\node[anchor=north east,inner sep=0pt,outer sep=0pt,align=left] at (199,268) {\fontsize{9}{11.34}\selectfont 5};
\end{scope}
\end{tikzpicture}
\newpage
\null
\begin{tikzpicture}[remember picture,overlay,x=1mm,y=-1mm]
\begin{scope}[shift={(current page.north west)}]
\node[anchor=north west,inner sep=0pt,outer sep=0pt,align=left,text width=187mm] at (15,12) {\fontsize{11}{13.86}\selectfont Week 26 / Same area, different boundaries / Grades 4--5};
\node[anchor=north west,inner sep=0pt,outer sep=0pt,align=left,text width=184mm] at (15,28) {\fontsize{14}{17.64}\selectfont \textbf{Problem 6:} Give a rule for the greatest number of squares in a shape with any even boundary length of at least 4. Explain why your rule works. Use it to find the shortest possible boundaries for shapes with 37, 50, and 73 squares, and draw shapes that reach them.};
\begin{scope}[shift={(24,92)}]
\draw[step=6mm,black!27,line width=.3pt] (0,0) grid (66,60);
\end{scope}
\draw[black!27,line width=.3pt] (24,92) rectangle (90,152);
\node[anchor=north west,inner sep=0pt,outer sep=0pt,align=left] at (24,155) {\fontsize{11}{13.86}\selectfont 37 squares; boundary: \rule{10mm}{.3pt}};
\begin{scope}[shift={(123,92)}]
\draw[step=6mm,black!27,line width=.3pt] (0,0) grid (66,60);
\end{scope}
\draw[black!27,line width=.3pt] (123,92) rectangle (189,152);
\node[anchor=north west,inner sep=0pt,outer sep=0pt,align=left] at (123,155) {\fontsize{11}{13.86}\selectfont 50 squares; boundary: \rule{10mm}{.3pt}};
\begin{scope}[shift={(24,181)}]
\draw[step=6mm,black!27,line width=.3pt] (0,0) grid (66,60);
\end{scope}
\draw[black!27,line width=.3pt] (24,181) rectangle (90,241);
\node[anchor=north west,inner sep=0pt,outer sep=0pt,align=left] at (24,244) {\fontsize{11}{13.86}\selectfont 73 squares; boundary: \rule{10mm}{.3pt}};
\draw[black!30,line width=.35pt] (113,186) -- ++(80,0);
\draw[black!30,line width=.35pt] (113,202) -- ++(80,0);
\draw[black!30,line width=.35pt] (113,218) -- ++(80,0);
\draw[black!30,line width=.35pt] (113,234) -- ++(80,0);
\draw[black!30,line width=.35pt] (113,250) -- ++(80,0);
\node[anchor=north west,inner sep=0pt,outer sep=0pt,align=left] at (15,268) {\fontsize{9}{11.34}\selectfont Bellingham Math Circle / Week 26 / W26-45-v2};
\node[anchor=north east,inner sep=0pt,outer sep=0pt,align=left] at (199,268) {\fontsize{9}{11.34}\selectfont 6};
\end{scope}
\end{tikzpicture}
\end{document}
